% Gemeinsame Vorlage für alle Übungsblätter.
% Aufgabenblatt:  pdflatex -jobname=UE01_Aufgaben UE01.tex
% Lösungsblatt:   pdflatex -jobname=UE01_Loesung "\def\mitloesung{}% Gemeinsame Vorlage für alle Übungsblätter.
% Aufgabenblatt:  pdflatex -jobname=UE01_Aufgaben UE01.tex
% Lösungsblatt:   pdflatex -jobname=UE01_Loesung "\def\mitloesung{}\input{praeambel}
\begin{document}
\blatt{1}{Grundlagen in R}{1}

\textbf{Vorlage:} \texttt{Exercises/UE01\_Aufgaben.R}\\
Schreiben Sie Ihren Code unter die jeweilige Aufgabe. Die Kontrollwerte zeigen, ob Ihr Ergebnis stimmt. R zeigt Dezimalzahlen mit Punkt (3.75), auf diesem Blatt stehen sie mit Komma (3,75).

\aufgabe{Rechnen}
Berechnen Sie $5 + 12$, $8 \cdot 4$, $15 / 4$ und $(7 \cdot 3) + (9 / 2)$.
\kontrolle{17; 32; 3,75; 25,5}

\aufgabe{Objekte}
Weisen Sie \texttt{a} den Wert 10 zu und \texttt{b} den Wert \texttt{a * 2}. Berechnen Sie \texttt{a + b}.
\kontrolle{30}

\aufgabe{Vergleiche}
Prüfen Sie mit R: $10 > 5$, $4 \le 8$, $3 \ne 2$ und $7 = 7$.
\kontrolle{viermal \texttt{TRUE}}
\deutung{R schreibt dafür \texttt{>}, \texttt{<=}, \texttt{!=} und \texttt{==}. Ein einfaches \texttt{=} weist einen Wert zu und vergleicht nicht.}

\aufgabe{Vektoren}
Erzeugen Sie \texttt{v <- c(3, 8, 1, 12, 7)}. Berechnen Sie Summe und Mittelwert, geben Sie das vierte Element aus und alle Werte größer als 5.
\kontrolle{31; 6,2; 12; 8 12 7}

\aufgabe{Fehlende Werte}
Erzeugen Sie \texttt{w <- c(2, NA, 6, 10)}. Berechnen Sie den Mittelwert einmal ohne und einmal mit \texttt{na.rm = TRUE}. Zählen Sie die fehlenden Werte mit \texttt{is.na()}.
\kontrolle{NA; 6; 1}
\deutung{Ein einziger fehlender Wert macht den Mittelwert unbekannt. Mit \texttt{na.rm = TRUE} bleibt er unberücksichtigt.}

\aufgabe{Datentypen}
Erzeugen Sie \texttt{x <- c(1, 2, "drei")}. Welche Klasse hat \texttt{x}? Was liefert \texttt{as.numeric(x)}? Erklären Sie das Ergebnis in einem Satz.
\kontrolle{\texttt{"character"}; 1 2 NA mit Warnung}
\deutung{Ein Vektor enthält nur einen Datentyp. Wegen \enquote{drei} wird alles zu Text, beim Zurückwandeln wird \enquote{drei} zu NA.}

\loesungscode{../UE01_Loesung.R}
\end{document}
"
\documentclass[11pt,a4paper]{article}
\usepackage[T1]{fontenc}
\usepackage[utf8]{inputenc}
\usepackage[ngerman]{babel}
\usepackage{lmodern}
\usepackage[a4paper,margin=2.2cm,top=2.5cm]{geometry}
\usepackage{xcolor}
\usepackage{listings}
\usepackage{fancyhdr}
\usepackage[autostyle]{csquotes}
\usepackage{parskip}
\usepackage{enumitem}

\AtBeginDocument{\shorthandoff{"}}

\newif\ifloesung
\ifdefined\mitloesung \loesungtrue \fi

\pagestyle{fancy}
\fancyhf{}
\lhead{\small Computergestützte Statistische Datenanalyse}
\rhead{\small HWR Berlin \textperiodcentered{} WiSe 2026}
\cfoot{\small\thepage}

\newcommand{\blatt}[3]{%
  {\LARGE\bfseries \ifloesung Lösung zu \fi Übungsblatt #1}\par\smallskip
  {\large #2}\par
  {\small zu Vorlesung #3}\par\medskip}

\newcounter{aufgabe}
\newcommand{\aufgabe}[1]{\stepcounter{aufgabe}\par\medskip\noindent\textbf{Aufgabe \theaufgabe: #1}\par\nopagebreak}
\newcommand{\kontrolle}[1]{\par\noindent{\small\itshape Kontrollwert: #1}\par}
\newcommand{\deutung}[1]{\ifloesung\par\noindent{\small\textbf{Musterdeutung:} #1}\par\fi}
\newcommand{\hinweis}[1]{\par\noindent{\small\textbf{Tipp:} #1}\par}

\lstset{language=R, basicstyle=\ttfamily\footnotesize, commentstyle=\color{gray}\itshape,
  keywordstyle=\color{blue!60!black}, stringstyle=\color{green!40!black},
  frame=single, rulecolor=\color{gray!40}, breaklines=true, columns=fullflexible,
  keepspaces=true, showstringspaces=false,
  literate={ä}{{\"a}}1 {ö}{{\"o}}1 {ü}{{\"u}}1 {Ä}{{\"A}}1 {Ö}{{\"O}}1 {Ü}{{\"U}}1
           {ß}{{\ss}}1 {²}{{\textsuperscript{2}}}1 {³}{{\textsuperscript{3}}}1}

\newcommand{\loesungscode}[1]{\ifloesung\section*{R-Code der Lösung}\lstinputlisting{#1}\fi}

\begin{document}
\blatt{1}{Grundlagen in R}{1}

\textbf{Vorlage:} \texttt{Exercises/UE01\_Aufgaben.R}\\
Schreiben Sie Ihren Code unter die jeweilige Aufgabe. Die Kontrollwerte zeigen, ob Ihr Ergebnis stimmt. R zeigt Dezimalzahlen mit Punkt (3.75), auf diesem Blatt stehen sie mit Komma (3,75).

\aufgabe{Rechnen}
Berechnen Sie $5 + 12$, $8 \cdot 4$, $15 / 4$ und $(7 \cdot 3) + (9 / 2)$.
\kontrolle{17; 32; 3,75; 25,5}

\aufgabe{Objekte}
Weisen Sie \texttt{a} den Wert 10 zu und \texttt{b} den Wert \texttt{a * 2}. Berechnen Sie \texttt{a + b}.
\kontrolle{30}

\aufgabe{Vergleiche}
Prüfen Sie mit R: $10 > 5$, $4 \le 8$, $3 \ne 2$ und $7 = 7$.
\kontrolle{viermal \texttt{TRUE}}
\deutung{R schreibt dafür \texttt{>}, \texttt{<=}, \texttt{!=} und \texttt{==}. Ein einfaches \texttt{=} weist einen Wert zu und vergleicht nicht.}

\aufgabe{Vektoren}
Erzeugen Sie \texttt{v <- c(3, 8, 1, 12, 7)}. Berechnen Sie Summe und Mittelwert, geben Sie das vierte Element aus und alle Werte größer als 5.
\kontrolle{31; 6,2; 12; 8 12 7}

\aufgabe{Fehlende Werte}
Erzeugen Sie \texttt{w <- c(2, NA, 6, 10)}. Berechnen Sie den Mittelwert einmal ohne und einmal mit \texttt{na.rm = TRUE}. Zählen Sie die fehlenden Werte mit \texttt{is.na()}.
\kontrolle{NA; 6; 1}
\deutung{Ein einziger fehlender Wert macht den Mittelwert unbekannt. Mit \texttt{na.rm = TRUE} bleibt er unberücksichtigt.}

\aufgabe{Datentypen}
Erzeugen Sie \texttt{x <- c(1, 2, "drei")}. Welche Klasse hat \texttt{x}? Was liefert \texttt{as.numeric(x)}? Erklären Sie das Ergebnis in einem Satz.
\kontrolle{\texttt{"character"}; 1 2 NA mit Warnung}
\deutung{Ein Vektor enthält nur einen Datentyp. Wegen \enquote{drei} wird alles zu Text, beim Zurückwandeln wird \enquote{drei} zu NA.}

\loesungscode{../UE01_Loesung.R}
\end{document}
"
\documentclass[11pt,a4paper]{article}
\usepackage[T1]{fontenc}
\usepackage[utf8]{inputenc}
\usepackage[ngerman]{babel}
\usepackage{lmodern}
\usepackage[a4paper,margin=2.2cm,top=2.5cm]{geometry}
\usepackage{xcolor}
\usepackage{listings}
\usepackage{fancyhdr}
\usepackage[autostyle]{csquotes}
\usepackage{parskip}
\usepackage{enumitem}

\AtBeginDocument{\shorthandoff{"}}

\newif\ifloesung
\ifdefined\mitloesung \loesungtrue \fi

\pagestyle{fancy}
\fancyhf{}
\lhead{\small Computergestützte Statistische Datenanalyse}
\rhead{\small HWR Berlin \textperiodcentered{} WiSe 2026}
\cfoot{\small\thepage}

\newcommand{\blatt}[3]{%
  {\LARGE\bfseries \ifloesung Lösung zu \fi Übungsblatt #1}\par\smallskip
  {\large #2}\par
  {\small zu Vorlesung #3}\par\medskip}

\newcounter{aufgabe}
\newcommand{\aufgabe}[1]{\stepcounter{aufgabe}\par\medskip\noindent\textbf{Aufgabe \theaufgabe: #1}\par\nopagebreak}
\newcommand{\kontrolle}[1]{\par\noindent{\small\itshape Kontrollwert: #1}\par}
\newcommand{\deutung}[1]{\ifloesung\par\noindent{\small\textbf{Musterdeutung:} #1}\par\fi}
\newcommand{\hinweis}[1]{\par\noindent{\small\textbf{Tipp:} #1}\par}

\lstset{language=R, basicstyle=\ttfamily\footnotesize, commentstyle=\color{gray}\itshape,
  keywordstyle=\color{blue!60!black}, stringstyle=\color{green!40!black},
  frame=single, rulecolor=\color{gray!40}, breaklines=true, columns=fullflexible,
  keepspaces=true, showstringspaces=false,
  literate={ä}{{\"a}}1 {ö}{{\"o}}1 {ü}{{\"u}}1 {Ä}{{\"A}}1 {Ö}{{\"O}}1 {Ü}{{\"U}}1
           {ß}{{\ss}}1 {²}{{\textsuperscript{2}}}1 {³}{{\textsuperscript{3}}}1}

\newcommand{\loesungscode}[1]{\ifloesung\section*{R-Code der Lösung}\lstinputlisting{#1}\fi}
